%============================================================
% CH1 · P5 · PROBLEM 1 — BADGE (Codeforces 1020B)
%============================================================
\subsection{Problem 1.5.1: Badge}
\problemlabel{1.5.1}{ssec:ch1-badge}

\problemcard{1.5.1}{Badge}{Codeforces 1020B}{https://codeforces.com/problemset/problem/1020/B}{\exeasy\ \cf{1000}}{Follow the arrows until a vertex repeats: the answer is the entry point of the $\rho$; \bigO{n} by filling tails from the cycles}

\subsubsection*{1. Problem statement}

A teacher caught $n$ students (numbered $1..n$). She goes to some student $a$ and makes a hole in his badge. That student blames student $p_a$; the teacher goes there and makes a hole in $p_a$'s badge; he blames $p_{p_a}$, and so on. Eventually she comes to a student who \emph{already} has a hole; she makes a second hole in his badge and stops. For \textbf{every} possible first student $a = 1..n$, print who gets the second hole.

\textbf{Constraints:} $1 \le n \le 1000$, \; $1 \le p_i \le n$ (a student may blame himself).

\begin{table}[H]
\centering\small
\begin{tabular}{@{}lll@{}}
\toprule
\textbf{Input} & \textbf{Output} & \textbf{Explanation} \\
\midrule
\texttt{3} / \texttt{2 3 2} & \texttt{2 2 3} & from 1: $1,2,3,\mathbf{2}$; from 2: $2,3,\mathbf{2}$; from 3: $3,2,\mathbf{3}$ \\
\texttt{3} / \texttt{1 2 3} & \texttt{1 2 3} & everybody blames himself: the first student gets both holes \\
\bottomrule
\end{tabular}
\end{table}

\begin{diagrambox}
\begin{tikzpicture}
    \begin{scope}
        \node[gv] (s1) at (0,0) {1};
        \node[gv] (s2) at (1.6,0) {2};
        \node[gv] (s3) at (3.2,0) {3};
        \draw[gde] (s1) -- (s2);
        \draw[gde] (s2) to[bend left=35] (s3);
        \draw[gde] (s3) to[bend left=35] (s2);
        \node[font=\small\bfseries] at (1.6,1.2) {Sample 1: $p = [2,3,2]$};
    \end{scope}
    \begin{scope}[xshift=6.6cm]
        \foreach \v/\x in {1/0, 2/1.6, 3/3.2} {
            \node[gv] (t\v) at (\x,0) {\v};
            \draw[gde] (t\v) to[out=60, in=120, looseness=6] (t\v);
        }
        \node[font=\small\bfseries] at (1.6,1.2) {Sample 2: $p = [1,2,3]$};
    \end{scope}
\end{tikzpicture}
\figcap{Each student is a vertex with one arrow ``blames''. Sample 1 is one $\rho$: tail $1$, cycle $2 \leftrightarrows 3$. Sample 2 has three self-loops (cycles of length 1).}
\end{diagrambox}

\subsubsection*{2. The graph model}

\begin{modelbox}[Graph Model --- Badge]
\renewcommand{\arraystretch}{1.35}
\begin{tabularx}{\linewidth}{@{}>{\bfseries}l X@{}}
Vertices & the $n$ students \\
Edges & $i \to p_i$ (``$i$ blames $p_i$''); exactly one outgoing edge per vertex: a \textbf{functional graph}, stored as the array \texttt{p} itself \\
Weights & none \\
Question & for every start $a$: the first vertex visited twice by the walk $a, p_a, p_{p_a}, \dots$ = the \textbf{entry point} of the $\rho$ that starts at $a$ \\
Size & $|V| = |E| = n \le 1000$; one walk has at most $n+1$ steps, so even $n$ walks ($10^6$ steps) are cheap \\
\end{tabularx}
\end{modelbox}

\begin{diagrambox}
\begin{tikzpicture}[x=1.45cm]
    \foreach \v/\i/\s in {1/0/gvsrc, 2/1/gvdst, 3/2/gvseen}{
        \node[\s] (w\i) at (\i,0) {\v};
        \node[gdist, below=1mm of w\i] {hole 1};
    }
    \node[gvdst, dashed] (w3) at (3,0) {2};
    \node[gdist, below=1mm of w3, text=warningred] {hole 2!};
    \foreach \i [evaluate=\i as \j using int(\i+1)] in {0,...,2} \draw[gde] (w\i) -- (w\j);
    \draw[gback] (w3.north) to[bend right=35] node[lbl, above]{already has a hole} (w1.north);
    \node[note, text width=5.6cm, align=left, anchor=west] at (3.6,0) {Sample 1, $a = 1$. The teacher's route is the $\rho$ of vertex 1; the student with two holes is the \textbf{first vertex that repeats} = the entry point, here $2$.};
\end{tikzpicture}
\end{diagrambox}

\subsubsection*{3. Key observations}

\begin{itemize}[leftmargin=*]
    \item \textbf{The walk is forced}, so for a fixed $a$ the answer is simply ``simulate until a repeat''. The process always stops within $n+1$ visits (pigeonhole).
    \item \textbf{The answer is the entry point.} Tail vertices are visited once, then the walk enters the cycle at $e$ and goes around; the first vertex it sees again is $e$ (shape theorem, \S\ref{ssec:ch1-p5-theory}).
    \item \textbf{A cycle vertex is its own answer.} If $a$ lies on a cycle, the walk goes once around and returns to $a$ first: $\text{ans}[a] = a$.
    \item \textbf{A tail vertex inherits the answer of its successor.} If $a$ is not on a cycle, the walks from $a$ and from $p_a$ enter the cycle at the same vertex: $\text{ans}[a] = \text{ans}[p_a]$. This is what makes an \bigO{n} solution possible.
\end{itemize}

\begin{warningbox}[Not ``the vertex with the largest in-degree'']
In Sample 1 vertex $2$ has in-degree 2 and is the answer for $a = 1, 2$ --- but for $a = 3$ the answer is $3$ (in-degree 1). The answer depends on where you \emph{enter} the cycle, not on popularity.
\end{warningbox}

\subsubsection*{4. Approaches}

\begin{bruteforcebox}[Approach 1: Direct Simulation (already accepted here)]
For each $a$, use a fresh array \texttt{hole[]} and walk until the current student already has a hole.

\begin{lstlisting}[style=cpp]
#include <bits/stdc++.h>
using namespace std;

int main() {
    int n;
    cin >> n;
    vector<int> p(n + 1);
    for (int i = 1; i <= n; i++) cin >> p[i];
    for (int a = 1; a <= n; a++) {
        vector<bool> hole(n + 1, false);  // fresh badges for every experiment
        int v = a;
        while (!hole[v]) {                // first visit: one hole, move on
            hole[v] = true;
            v = p[v];
        }
        cout << v << (a == n ? '\n' : ' ');  // v is the first student visited twice
    }
}
\end{lstlisting}

\begin{complexitybox}
\textbf{Time:} \bigO{n^2} --- at most $n+1$ steps per start, about $10^6$ for $n = 1000$. \textbf{Space:} \bigO{n}. \textbf{Verdict:} \ac{} for this problem (it is the intended solution), but \tle{} if $n$ were $2\cdot10^5$ (about $2\cdot10^{10}$ steps on one long chain).
\end{complexitybox}
\end{bruteforcebox}

\begin{algorithmbox}[Approach 2: Linear Time --- Find the Cycles, Then Fill the Tails]
Use Algorithm 2 of the theory (three states). When a walk closes a \emph{new} cycle, every cycle vertex is its own answer. Then the tail vertices of this walk are answered \textbf{backwards}, from the cycle outwards, with $\text{ans}[v] = \text{ans}[p_v]$. When a walk instead runs into a finished vertex, that vertex already has its answer, and again the path is filled backwards.

\begin{lstlisting}[style=cpp]
#include <bits/stdc++.h>
using namespace std;

int main() {
    int n;
    cin >> n;
    vector<int> p(n + 1);
    for (int i = 1; i <= n; i++) cin >> p[i];

    vector<int> state(n + 1, 0);          // 0 = new, 1 = on the current walk, 2 = answered
    vector<int> ans(n + 1, 0);
    for (int s = 1; s <= n; s++) {
        if (state[s] != 0) continue;
        vector<int> path;
        int v = s;
        while (state[v] == 0) {           // walk until we meet something already seen
            state[v] = 1;
            path.push_back(v);
            v = p[v];
        }
        int k = path.size();              // path[0..k-1] still need an answer
        if (state[v] == 1) {              // NEW cycle at v: cycle vertices answer themselves
            while (path[k - 1] != v) { ans[path[k - 1]] = path[k - 1]; k--; }
            ans[v] = v; k--;
        }
        for (int i = k - 1; i >= 0; i--)  // tail vertices, from the cycle outwards
            ans[path[i]] = ans[p[path[i]]];
        for (int u : path) state[u] = 2;
    }
    for (int a = 1; a <= n; a++) cout << ans[a] << (a == n ? '\n' : ' ');
}
\end{lstlisting}

\begin{complexitybox}
\textbf{Time:} \bigO{n} --- every vertex enters a \texttt{path} once and is answered once. \textbf{Space:} \bigO{n}. \textbf{Verdict:} \ac{}, and it would still be fast for $n = 10^6$.
\end{complexitybox}
\end{algorithmbox}

\begin{invariantbox}[Why it is correct]
\textbf{Claim:} when a vertex is answered, its answer is the entry point of its $\rho$.

\emph{Cycle vertices.} The walk stopped at a state-1 vertex $v$, so $v, p_v, \dots$ up to the end of \texttt{path} is a cycle (theory, Algorithm 2). Starting on a cycle, the first vertex seen twice is the start itself: $\text{ans}[u] = u$ is right.

\emph{Tail vertices.} For a vertex $u$ not on a cycle, the walk from $u$ is $u$ followed by the walk from $p_u$; $u$ itself is never visited again (it is not on a cycle), so the first repeated vertex of both walks is the same: $\text{ans}[u] = \text{ans}[p_u]$. Filling \texttt{path} \emph{backwards} guarantees that $p_u$ (the next element of \texttt{path}, or the vertex $v$ where the walk stopped) is already answered when $u$ is processed. \hfill$\blacksquare$
\end{invariantbox}

\begin{alternativebox}[Approach 3: Peeling (Kahn) --- Cycles First, Then Reverse Order]
Delete in-degree-0 vertices repeatedly (theory, Algorithm 3). The survivors are the cycle vertices, $\text{ans}[v] = v$. A deleted vertex always comes out \emph{before} its successor, so processing the deletion order \textbf{backwards} sees $p_v$ before $v$, and $\text{ans}[v] = \text{ans}[p_v]$ can be filled directly.

\begin{lstlisting}[style=cpp]
#include <bits/stdc++.h>
using namespace std;

int main() {
    int n;
    cin >> n;
    vector<int> p(n + 1), indeg(n + 1, 0);
    for (int i = 1; i <= n; i++) { cin >> p[i]; indeg[p[i]]++; }

    vector<int> order;                    // vertices in the order they are peeled
    queue<int> q;
    for (int v = 1; v <= n; v++) if (indeg[v] == 0) q.push(v);
    while (!q.empty()) {
        int v = q.front(); q.pop();
        order.push_back(v);
        if (--indeg[p[v]] == 0) q.push(p[v]);
    }
    vector<int> ans(n + 1, 0);
    for (int v = 1; v <= n; v++)
        if (indeg[v] > 0) ans[v] = v;     // survivors lie on cycles
    for (int i = (int)order.size() - 1; i >= 0; i--)  // p[v] is peeled after v,
        ans[order[i]] = ans[p[order[i]]];             // so it is answered before v
    for (int a = 1; a <= n; a++) cout << ans[a] << (a == n ? '\n' : ' ');
}
\end{lstlisting}

\begin{diagrambox}
\begin{tikzpicture}
    \foreach \v/\x/\y/\s in {1/-1.4/0.8/gvdead, 2/0/0.8/gvsel, 3/1.3/0/gvsel, 4/0/-0.8/gvsel, 5/-1.4/-0.8/gvdead, 9/-2.8/-0.8/gvdead,
                             6/3.2/0.8/gvsel, 7/4.5/0/gvsel, 8/3.2/-0.8/gvsel}
        \node[\s] (k\v) at (\x,\y) {\v};
    \foreach \u/\w in {2/3, 3/4, 4/2, 6/7, 7/8, 8/6} \draw[gde] (k\u) -- (k\w);
    \foreach \u/\w in {1/2, 5/4, 9/5} \draw[gfaded, dashed, -{Stealth[length=2mm]}] (k\u) -- (k\w);
    \node[gdist, above=1mm of k1] {peeled 1st};
    \node[gdist, above=1mm of k9] {peeled 2nd};
    \node[gdist, below=1mm of k5] {peeled 3rd};
    \node[note, text width=4.4cm, align=left, anchor=west] at (5.2,0) {Peeling order $1, 9, 5$. Backwards: $\text{ans}[5] = \text{ans}[4] = 4$, then $\text{ans}[9] = \text{ans}[5] = 4$, then $\text{ans}[1] = \text{ans}[2] = 2$.};
\end{tikzpicture}
\figcap{Approach 3 on the trace input $p = [2,3,4,2,4,7,8,6,5]$ of Section 5. Purple: survivors = cycle vertices, each its own answer.}
\end{diagrambox}

\begin{complexitybox}
\textbf{Time:} \bigO{n}. \textbf{Space:} \bigO{n}. \textbf{Verdict:} \ac{}. Same answers, different order of work: Approach 2 discovers cycles while walking, Approach 3 first removes everything that is not on a cycle.
\end{complexitybox}
\end{alternativebox}

\begin{tricksbox}[Approach 4: Implementation Notes]
\begin{itemize}[leftmargin=*]
    \item \textbf{No recursion.} The tempting \texttt{solve(v)} that returns \texttt{ans[v] = solve(p[v])} recurses along a tail of length up to $n$; with $n = 1000$ that is safe, with $n = 10^6$ it overflows the stack. The explicit \texttt{path} vector is the iterative version of that recursion.
    \item \textbf{1-indexing.} Students are $1..n$; arrays of size $n+1$ avoid every $\pm 1$.
    \item \textbf{Output.} The $n$ answers go on one line separated by spaces; the code prints a newline after the last one instead of a space.
\end{itemize}
\end{tricksbox}

\subsubsection*{5. Trace}

We trace Approach 2 on a bigger input with two cycles, a tail of length 2 and two tails entering the same cycle at different vertices:\quad $n = 9$, \;$p = [2,3,4,2,4,7,8,6,5]$. The program prints \mbox{\texttt{2 2 3 4 4 6 7 8 4}}.

\begin{diagrambox}
\begin{tikzpicture}
    \foreach \v/\x/\y in {1/-1.4/0.8, 2/0/0.8, 3/1.3/0, 4/0/-0.8, 5/-1.4/-0.8, 9/-2.8/-0.8, 6/3.2/0.8, 7/4.5/0, 8/3.2/-0.8}
        \node[gv] (g\v) at (\x,\y) {\v};
    \foreach \u/\w in {1/2, 2/3, 3/4, 4/2, 5/4, 9/5, 6/7, 7/8, 8/6} \draw[gde] (g\u) -- (g\w);
\end{tikzpicture}
\figcap{The trace input: cycle $2 \to 3 \to 4 \to 2$ with tails $1 \to 2$ and $9 \to 5 \to 4$; cycle $6 \to 7 \to 8 \to 6$.}
\end{diagrambox}

\begin{tracebox}[{Approach 2 on $p = [2,3,4,2,4,7,8,6,5]$}]
\centering
\begin{tikzpicture}
    \begin{scope}
        \foreach \v/\x/\y/\s in {1/-1.4/0.8/gvseen, 2/0/0.8/gvseen, 3/1.3/0/gvseen, 4/0/-0.8/gvseen, 5/-1.4/-0.8/gv, 9/-2.8/-0.8/gv,
                                 6/3.2/0.8/gv, 7/4.5/0/gv, 8/3.2/-0.8/gv}
            \node[\s] (a\v) at (\x,\y) {\v};
        \foreach \u/\w in {1/2, 2/3, 3/4, 5/4, 9/5, 6/7, 7/8, 8/6} \draw[gde] (a\u) -- (a\w);
        \draw[gback] (a4) -- (a2);
        \node[note, text width=5.2cm, align=left, anchor=west] at (5.2,0) {\textbf{$s=1$:} \texttt{path} $= 1,2,3,4$, stops at $2$ (state 1). New cycle $2,3,4$: they are their own answers; then the tail: $\text{ans}[1] = \text{ans}[2] = 2$.};
    \end{scope}
    \begin{scope}[yshift=-2.75cm]
        \foreach \v/\x/\y/\s in {1/-1.4/0.8/gvdone, 2/0/0.8/gvdone, 3/1.3/0/gvdone, 4/0/-0.8/gvdone, 5/-1.4/-0.8/gvseen, 9/-2.8/-0.8/gv,
                                 6/3.2/0.8/gv, 7/4.5/0/gv, 8/3.2/-0.8/gv}
            \node[\s] (b\v) at (\x,\y) {\v};
        \foreach \u/\w in {1/2, 2/3, 3/4, 4/2, 9/5, 6/7, 7/8, 8/6} \draw[gde] (b\u) -- (b\w);
        \draw[gcross] (b5) -- (b4);
        \node[note, text width=5.2cm, align=left, anchor=west] at (5.2,0) {\textbf{$s=5$} ($2,3,4$ skipped: finished): \texttt{path} $= 5$, stops at $4$ (state 2, answered). No new cycle: $\text{ans}[5] = \text{ans}[4] = 4$.};
    \end{scope}
    \begin{scope}[yshift=-5.5cm]
        \foreach \v/\x/\y/\s in {1/-1.4/0.8/gvdone, 2/0/0.8/gvdone, 3/1.3/0/gvdone, 4/0/-0.8/gvdone, 5/-1.4/-0.8/gvdone, 9/-2.8/-0.8/gv,
                                 6/3.2/0.8/gvseen, 7/4.5/0/gvseen, 8/3.2/-0.8/gvseen}
            \node[\s] (c\v) at (\x,\y) {\v};
        \foreach \u/\w in {1/2, 2/3, 3/4, 4/2, 5/4, 9/5, 6/7, 7/8} \draw[gde] (c\u) -- (c\w);
        \draw[gback] (c8) -- (c6);
        \node[note, text width=5.2cm, align=left, anchor=west] at (5.2,0) {\textbf{$s=6$:} \texttt{path} $= 6,7,8$, stops at $6$ (state 1). A new cycle with no tail: $\text{ans}[8]=8$, $\text{ans}[7]=7$, $\text{ans}[6]=6$.};
    \end{scope}
    \begin{scope}[yshift=-8.25cm]
        \foreach \v/\x/\y/\s in {1/-1.4/0.8/gvdone, 2/0/0.8/gvdone, 3/1.3/0/gvdone, 4/0/-0.8/gvdone, 5/-1.4/-0.8/gvdone, 9/-2.8/-0.8/gvseen,
                                 6/3.2/0.8/gvdone, 7/4.5/0/gvdone, 8/3.2/-0.8/gvdone}
            \node[\s] (d\v) at (\x,\y) {\v};
        \foreach \u/\w in {1/2, 2/3, 3/4, 4/2, 5/4, 6/7, 7/8, 8/6} \draw[gde] (d\u) -- (d\w);
        \draw[gcross] (d9) -- (d5);
        \node[note, text width=5.2cm, align=left, anchor=west] at (5.2,0) {\textbf{$s=9$} ($7, 8$ skipped): \texttt{path} $= 9$, stops at $5$ (state 2): $\text{ans}[9] = \text{ans}[5] = 4$. Done.};
    \end{scope}
\end{tikzpicture}

\medskip
\small
\begin{tabular}{@{}c|ccccccccc@{}}
\toprule
$a$ & 1 & 2 & 3 & 4 & 5 & 6 & 7 & 8 & 9 \\
$p_a$ & 2 & 3 & 4 & 2 & 4 & 7 & 8 & 6 & 5 \\
\midrule
answered in walk & $s{=}1$ & $s{=}1$ & $s{=}1$ & $s{=}1$ & $s{=}5$ & $s{=}6$ & $s{=}6$ & $s{=}6$ & $s{=}9$ \\
\textbf{ans}$[a]$ & \textbf{2} & \textbf{2} & \textbf{3} & \textbf{4} & \textbf{4} & \textbf{6} & \textbf{7} & \textbf{8} & \textbf{4} \\
\bottomrule
\end{tabular}
\end{tracebox}

\begin{diagrambox}
\begin{tikzpicture}
    \foreach \v/\x/\y/\s/\a in {1/-1.4/0.8/gv/2, 2/0/0.8/gvdst/2, 3/1.3/0/gvdst/3, 4/0/-0.8/gvdst/4, 5/-1.4/-0.8/gv/4, 9/-2.8/-0.8/gv/4,
                                6/3.2/0.8/gvdst/6, 7/4.5/0/gvdst/7, 8/3.2/-0.8/gvdst/8}{
        \node[\s] (z\v) at (\x,\y) {\v};
        \node[gdist, text=accentorange, above right=-1.5mm and -1mm of z\v] {\a};
    }
    \foreach \u/\w in {1/2, 2/3, 3/4, 4/2, 5/4, 9/5, 6/7, 7/8, 8/6} \draw[gde] (z\u) -- (z\w);
    \node[note, text width=4.6cm, align=left, anchor=west] at (5.4,0) {Final answers (small orange numbers). Cycle vertices (orange) answer themselves; the tail $9 \to 5$ inherits $4$, the tail $1$ inherits $2$: the two tails enter the same cycle at different points.};
\end{tikzpicture}
\end{diagrambox}

Check $a = 9$ by hand: the teacher visits $9, 5, 4, 2, 3$ and then $4$ again --- student $4$ gets two holes, although the walk from $1$ enters the same cycle at $2$.

\subsubsection*{6. Test suite}

\begin{testbox}
\centering\small
\begin{tabular}{@{}lll@{}}
\toprule
\textbf{Input ($n$ / $p$)} & \textbf{Expected} & \textbf{Why this test} \\
\midrule
\texttt{3 / 2 3 2} & \texttt{2 2 3} & official sample 1: tail + 2-cycle \\
\texttt{3 / 1 2 3} & \texttt{1 2 3} & official sample 2: only self-loops \\
\texttt{1 / 1} & \texttt{1} & smallest input \\
\texttt{9 / 2 3 4 2 4 7 8 6 5} & \texttt{2 2 3 4 4 6 7 8 4} & two pieces, tail of length 2 (the trace) \\
\texttt{4 / 2 3 4 4} & \texttt{4 4 4 4} & a chain ending in a self-loop \\
$n = 1000$, $p_i = i+1$, $p_n = n$ & $n$ times \texttt{1000} & longest tail: brute force does $\approx 5\cdot10^5$ steps \\
\bottomrule
\end{tabular}
\end{testbox}

All three programs were also compared on 500 random inputs with $n \le 9$: identical outputs.

\subsubsection*{7. Common mistakes}

\begin{warningbox}[Mistake 1 --- forgetting to reset the holes between experiments]
In Approach 1 the array \texttt{hole[]} must be fresh for every $a$ (each $a$ is a separate story). Reusing it makes the second experiment stop immediately at a student who got his hole in the first one.
\end{warningbox}

\begin{warningbox}[Mistake 2 --- answering with the last new student instead of the repeated one]
\begin{lstlisting}[style=snippet]
while (!hole[p[v]]) { hole[v] = true; v = p[v]; }   // WRONG: stops one student too early
while (!hole[v]) { hole[v] = true; v = p[v]; }      // RIGHT: stop ON the repeated student
\end{lstlisting}
The answer is the student who receives the \emph{second} hole --- the vertex where the walk stops, not the one before it.
\end{warningbox}

\begin{warningbox}[Mistake 3 --- giving every vertex of the path the entry point]
In Approach 2, setting $\text{ans}[u] = v$ for the whole \texttt{path} is wrong for cycle vertices: on the cycle each vertex is its own answer ($\text{ans}[3] = 3$, not $2$). Only tail vertices inherit.
\end{warningbox}

\subsubsection*{8. Takeaways}

\begin{itemize}[leftmargin=*]
    \item ``Follow until something repeats'' $\Rightarrow$ functional graph; the first repeated vertex is the \textbf{entry point} of the $\rho$.
    \item Answers per vertex: cycle vertices are special (here: their own answer), tail vertices \textbf{inherit from their successor}.
    \item Fill tails backwards along the walk (or in reverse peeling order) so that the successor is always answered first.
    \item Small $n$: plain simulation is a perfectly good solution --- but know the \bigO{n} version for when $n$ grows.
\end{itemize}

\begin{variantsbox}
\variant{Longest cycle}{\varin{https://leetcode.com/problems/longest-cycle-in-a-graph/}}
  {instead of an entry point per start, report the length of the longest cycle; some vertices have no arrow.}
  {the same functional graph with \texttt{-1} for missing arrows; only cycles matter, tails are ignored.}
  {three-state walks (or visit times); a walk that stops inside itself gives a cycle length; keep the maximum.}
\variant{Where after $K$ visits?}{\varin{https://atcoder.jp/contests/abc167/tasks/abc167_d}}
  {one start only, but which vertex is reached after exactly $K \le 10^{18}$ steps?}
  {the $\rho$ of the start: tail $\mu$, cycle $\lambda$, entry point as here.}
  {walk until the first repeat, then $K \to \mu + (K - \mu) \bmod \lambda$ when $K \ge \mu$.}
\variant{Number of holed students}{\varnone}
  {for every $a$, how many \emph{different} students get at least one hole?}
  {per vertex: distance $d[a]$ to its cycle and length $\lambda[a]$ of that cycle; the answer is $d[a] + \lambda[a] = \mu + \lambda$.}
  {cycle vertices: $d = 0$; tails filled from the cycle outwards: $d[a] = d[p_a] + 1$, $\lambda[a] = \lambda[p_a]$ --- the skeleton of Approach 2 (more in Chapter~11).}
\variant{Trap the student cheaply}{\varin{https://codeforces.com/problemset/problem/1027/D}}
  {each vertex has a price; choose vertices of minimum total price so that \emph{every} walk eventually hits a chosen one.}
  {each walk ends in its cycle, so every cycle needs one chosen vertex; tails need none.}
  {for every new cycle found by the three-state walk, add its cheapest vertex.}
\end{variantsbox}
